\section{Modeling Finite Precision Arithmetic}

To bridge theory and practice, we want to formulate our results on a structure that exists under computer arithmetic constraints. Modeling classifiers over real numbers can lead to a mismatch between theoretical insights and practical implementations. For example, verifiers for neural networks based on real arithmetic can be unsound when applied to floating-point computations~\cite{jia2021exploiting,cordeiro2025neural}.

\paragraph{Mathematcal Model.} Floating point numbers approximate real values with finite precision according to the IEEE 754 standard, which defines a finite arithmetic over a set of bitvector-encoded numbers.
Reasoning about finite-precision arithmetic with standard operators (such as addition) is challenging, since these operators may violate familiar algebraic properties like associativity~\cite{Kneusel2025}. 
An additional problem is presented by special floating point values such as \texttt{NaN} or \texttt{Inf}.
Nevertheless, the IEEE 754-2008 standard provides a well-defined total ordering on floating-point numbers~\cite{4610935}, implemented, for example, in Rust\footnote{\url{https://doc.rust-lang.org/std/primitive.f64.html\#method.total\_cmp}}. This motivates formulation of our results over ordered spaces $(S, <, d)$.

But an introduction of an order will not lead to more productive results unless a relationship between the order and the distance exists. To fix this, we introduce a new inequality. The inequality states that visiting an intermediate point between two points in the order produces at most the direct distance. In some sense, it is the opposite of a triangle inequality, which states that visiting an additional point increases the distance. But as our goal in further proofs is to bound the size of balls, we found this inequality more productive. We also decided against making this an equality for a irrelevant to our further results possibility that distance uses floating point rounding.

With these considerations, we arrive at our mathematical model standing for computer arithmetic spaces:
% lax begin Lax544225.OrderedSemimetric
\begin{definition}[Ordered Semimetric Space]
    \label{def:ordered-space}
    A tuple $(S, <, d)$ with a set $S$, a (partial) order $< \subseteq S \times S$, and $d: S \times S \rightarrow \mathbb{R}_{\geq 0}$ is an \emph{ordered semimetric space} if $(S, d)$ is a semimetric space and for all $a,b,c \in S$ such that $a \leq b \leq c$ implies $d(a, c) \geq d(a, b) + d(b,c)$.
\end{definition}
% lax end


\paragraph{Model Validation.} After defining ordered semimetric spaces, the question remains whether such spaces indeed exist on floating point numbers and can resemble common distances. To verify this, we utilized Mistral Vibe to generate provably correct instances of ordered semimetric spaces with respect to the IEEE 754 total order on lists of Flocq~\cite{boldo2011flocq} floating point numbers within the Rocq proof asisstant~\cite{rocq}, re-implement the resulting distances for Python  64 bit \texttt{float} and produce a script for differential testing between Rocq and Python using fixed edge cases and randomly generated instances. The task was successfully done by Mistral Vibe in 6 hours and did cost approximately 45 euro. The output was reviewed by humans afterwards and is avaliable on GitHub\footnote{\url{https://github.com/verinncoq/mueps-robustness-proofs/tree/main/flocq_ordered_spaces}}. The following proposition summarizes the findings.

\begin{proposition} 
\label{prop:fp-space}
% lax begin Lax544225.BinaryFloats
Let $\mathcal{F}$ be a binary floating point format with $p \in \mathbb{N}$ being the number of mantissa bits. Then, it consists of triples $(s, c, m)$ representing numbers $(-1)^s \cdot 2^c \cdot m$ such that $s \in \{0,1\}$, $c \in \{c_{\min},\dots, c_{\max}\} \subset \mathbb{Z}$, $m \in \{k \cdot 2^{-p} : k \in \{1, \dots, 2^{p+1}-1\} \subset \mathbb{Z}\}$ together with additional values $\pm0$, \texttt{NaN}, $\pm$\texttt{Inf} $\in \mathcal{F}$. For each $(s, c, m) \in \mathcal{F}$ holds that either $m \geq 1$ (normal) or $c = c_{\min}$ (subnormal). Let $\phi : \mathcal{F} \rightarrow \mathbb{R}$ with:
\begin{itemize}
    \item $\phi((s,c,m)) = (-1)^s \cdot 2^c \cdot m$ 
    \item $\phi(+0) = 0$ and $\phi(-0) = -2^{c_{\min} - p - 1}$
    \item $\phi(\texttt{NaN}) = 2^{c_{\max} + 1}  + 1$
    \item $\phi(\pm\texttt{Inf}) = \pm 2^{c_{\max} + 1}$
\end{itemize}
% lax end
Then, the following are ordered semimetric spaces:
\begin{enumerate}
% lax begin Lax544225.FloatAbsoluteSpace
    \item $(\mathcal{F}, <_{IEEE}, d_1)$ where $<_{IEEE}$ is the IEEE total order and $d_1(f_1, f_2) = |\phi(f_1) - \phi(f_2)|$.
% lax end
% lax begin Lax544225.FloatSquaredSpace
    \item $(\mathcal{F}, <_{IEEE}, d_2)$ where $d_2(f_1, f_2) = (\phi(f_1) - \phi(f_2))^2$.
% lax end
% lax begin Lax544225.ProductSpaces
    \item $(\mathcal{F}^n, <^n, d_3(f, g) := \sum_{i \in [n]} d(f_i, g_i))$
% lax end
% lax begin Lax544225.SumProductSpace
 such that $(\mathcal{F}, <, d)$ is an ordered semimetric space and $<^n$ is n-wise product partial order.
% lax end
% lax begin Lax544225.MaximumProductSpace
    \item $(\mathcal{F}^n, <^{n}_*, d_4(f, g) := \max_{i \in [n]}d(f_i, g_i))$ where $(\mathcal{F},<,d)$ is an ordered semimetric space such that for all $a,b,c\in\mathcal{F}$ with $a \leq b \leq c$ holds that $d(a,c) = d(a,b) + d(b,c)$ and $<^n_* := \{(f,g) \in \mathcal{F}^n \times \mathcal{F}^n\,|f <^n g \wedge \forall i,j \in [n]\, d(f_i, g_i) = d(f_j, g_j)\,\}$.
% lax end
    \item Extensions of
% lax begin Lax544225.SumListSpace
(3)
% lax end
and
% lax begin Lax544225.MaximumListSpace
(4)
% lax end
% lax begin Lax544225.ListSpaces
to arbitrary finite lists of floating-point numbers where a shorter list is padded with \texttt{NaN} for distance computation, and a penalty equal to $1$ is added to the distance whenever list lengths are different. The partial orders compare only lists of the same length, and lists of different lengths are incomparable.
% lax end
\end{enumerate}
\end{proposition}

The combinations of spaces 1 and 3, 1 and 4 in Proposition~\ref{prop:fp-space} give us notions similar to $l_1$, $l_\infty$ distances respectively. Overall, we showed that ordered semimetric spaces can be faithfully constructed on floating point numbers, showing that the results proven in the following sections are transferrable to computer arithmetic.